\section{Nonlinear polynomial fixed points, compressed packet composition, and finite tree signatures}
\label{sec:block11-nonlinear-tree}

Nonlinear tree atlas develops the nonlinear and rooted-tree stratum after the algebraic milestone.  The existing nonlinear sources already prove exact packet jets, the packet Fa\`a di Bruno law, colored operadic composition, refinement coherence, rooted-tree reconstruction, analytic majorant estimates, and local near-identity inversion \cite{P3,RT,NLO}.  The remaining problem is not whether nonlinear maps can be expanded.  It is whether a declared nonlinear family admits a normalized finite presentation and a finite signature that controls a named packet complexity.

This section gives three restricted answers.  First, a marked polynomial fixed-point system is organized into equivalent system, packet-hypergraph, operadic-tree, and exact packet-equation charts.  Second, repeated-direction packet Fa\`a di Bruno sums admit an exact orbit compression indexed by colored integer partitions rather than all labeled set partitions.  Third, simply generated plane-tree equations possess a finite arity-gcd signature that determines every packet-support pattern and, under the classical smooth implicit schema, the complete dominant-singularity orbit after arithmetic packetization.  A dependency condensation isolates the exact cyclic nonlinear core of a marked packet system.  None of these results is promoted to a universal nonlinear normal form.

Throughout the analytic part, $X$ is a Banach space over $\mathbb R$ or $\mathbb C$.  A finite packet system on $X$ is a family of bounded projections
\[
 P_iP_j=\delta_{ij}P_i,
 \qquad
 \sum_{i\in I}P_i=I_X.
\]
All sums over packet addresses are finite.  The formal-series part is valid over any characteristic-zero field unless positivity or singularity analysis is explicitly invoked.

\subsection{Marked polynomial fixed-point systems}
\label{subsec:block11-marked-systems}

Let
\[
 A_m:X^m\longrightarrow X,
 \qquad 2\le m\le D,
\]
be bounded symmetric multilinear maps, and let $s\in X$.  Consider the marked equation
\[
 \boxed{
 u=s+\Phi(u),
 \qquad
 \Phi(u)=\sum_{m=2}^{D}A_m[u,\ldots,u].}
 \tag{11.1}
\]
The word \emph{marked} means that the source term, every multilinear operation, the packet system, and the selected convergence or filtration certificate are retained as part of the presentation.  The fixed point $u$ alone is not assumed to reconstruct those data.

Put
\[
 a=\lVert s\rVert,
 \qquad
 b_m=\lVert A_m\rVert,
 \qquad
 \varphi(t)=\sum_{m=2}^{D}b_mt^m.
\]
For a radius $R>0$, define
\[
 \kappa_R=\varphi'(R)
 =\sum_{m=2}^{D}m b_mR^{m-1}.
\]

\begin{theorem}[Selected polynomial fixed point and tree expansion]
\label{thm:block11-polynomial-fixed-point}
Assume
\[
 a+\varphi(R)\le R,
 \qquad
 \kappa_R<1.
\]
Then equation \textup{(11.1)} has a unique solution in the closed ball $\overline B_R$.  Moreover, that solution is the absolutely convergent sum
\[
 \boxed{
 u=\sum_{T\in\mathcal T_{\mathcal A}}\operatorname{ev}(T),}
 \tag{11.2}
\]
where $\mathcal A=\{m:A_m\neq0\}$, the leaves of $T$ evaluate to $s$, and an internal vertex of arity $m$ evaluates by $A_m$ applied to the ordered values of its children.
\end{theorem}

\begin{proof}
The map
\[
 \Psi(x)=s+\Phi(x)
\]
preserves $\overline B_R$ because
\[
 \lVert\Psi(x)\rVert
 \le a+\sum_{m=2}^{D}b_mR^m
 \le R.
\]
For $x,y\in\overline B_R$, multilinearity gives
\[
 A_m[x^m]-A_m[y^m]
 =
 \sum_{r=1}^{m}
 A_m[y^{r-1},x-y,x^{m-r}],
\]
so
\[
 \lVert\Psi(x)-\Psi(y)\rVert
 \le
 \kappa_R\lVert x-y\rVert.
\]
The contraction theorem gives the unique fixed point.

Let $u^{(0)}=0$ and $u^{(h+1)}=\Psi(u^{(h)})$.  Expanding recursively shows that $u^{(h)}$ is the sum of the evaluations of all allowed plane rooted trees of height at most $h-1$.  If $r_0=0$ and
\[
 r_{h+1}=a+\varphi(r_h),
\]
then the sum of the norms of those tree evaluations is at most $r_h$.  The scalar iterates increase and remain bounded by $R$, hence converge.  Monotone convergence of the nonnegative norm majorant proves absolute convergence of the full tree sum.  Regrouping trees by the arity of the root gives equation \textup{(11.1)}.
\end{proof}

\begin{remark}[Formal filtered version]
If $X$ is replaced by a complete filtered module, $s$ begins in positive filtration, and every $A_m$ raises total filtration, then the same tree expansion is coefficientwise finite at each filtration depth.  It therefore defines the unique formal solution without a Banach norm.  The analytic and formal doctrines are related but are not silently identified.
\end{remark}

\begin{boundary}[The selected object is a marked system]
Two different polynomial systems can have the same fixed point.  Conversely, a single polynomial system can have several fixed points outside the contraction ball.  The normalized object in this section is therefore
\[
 (s,A_2,\ldots,A_D;\mathcal P;\text{certificate}),
\]
not an unmarked semantic vector $u$.
\end{boundary}

\subsection{Packet jets, hypergraphs, and the marked nonlinear atlas}
\label{subsec:block11-packet-atlas}

For $i,j_1,\ldots,j_m\in I$, define the resolved multilinear block
\[
 A_{m;i\mid j_1,\ldots,j_m}
 =
 P_iA_m(P_{j_1}\,\cdot,\ldots,P_{j_m}\,\cdot).
\]
Completeness of the packet family gives the exact reconstruction
\[
 \boxed{
 A_m
 =
 \sum_{i,j_1,\ldots,j_m}
 A_{m;i\mid j_1,\ldots,j_m}.}
 \tag{11.3}
\]
Write
\[
 u_i=P_iu,
 \qquad
 s_i=P_is.
\]
Then equation \textup{(11.1)} is equivalent to the complete packet system
\[
 \boxed{
 u_i
 =
 s_i+
 \sum_{m=2}^{D}
 \sum_{j_1,\ldots,j_m}
 A_{m;i\mid j_1,\ldots,j_m}
 [u_{j_1},\ldots,u_{j_m}].}
 \tag{11.4}
\]

The active nonlinear transfer hypergraph has packet addresses as colors and one typed hyperedge
\[
 (j_1,\ldots,j_m)\longrightarrow i
\]
for every nonzero block $A_{m;i\mid j_1,\ldots,j_m}$.  Rooted substitution of these hyperedges is precisely the colored symmetric operad already present in the nonlinear core sources \cite{RT,NLO}.

\begin{theorem}[Normalized atlas of a marked polynomial system]
\label{thm:block11-marked-atlas}
Fix a finite packet system and a certificate satisfying the hypotheses of \cref{thm:block11-polynomial-fixed-point}, or a declared complete filtration satisfying the formal version.  The following presentations determine one another exactly:
\[
 \boxed{
 \begin{aligned}
 &\text{the marked polynomial system }(s,A_2,\ldots,A_D),\\
 &\text{the complete family of packet blocks }A_{m;i\mid j_1,\ldots,j_m},\\
 &\text{the colored operadic tree grammar with its actual vertex operations},\\
 &\text{the exact packet fixed-point system \textup{(11.4)}}.
 \end{aligned}}
\]
Each chart reconstructs the same selected fixed point.  With literal normalized data and equality morphisms, every semantic fiber is contractible.
\end{theorem}

\begin{proof}
Equation \textup{(11.3)} reconstructs each multilinear operation from its complete block family.  The active hypergraph together with the actual block maps is simply another indexing of the same data.  Rooted operadic substitution generates exactly the tree evaluations of \cref{thm:block11-polynomial-fixed-point}.  Applying the output projections to the fixed-point equation gives \textup{(11.4)}, and synthesis recovers the original equation.  Once the packet system, block maps, source term, and certificate are normalized literally, an automorphism fixing the semantic object and all marked data is the identity.
\end{proof}

\begin{figure}[htbp]
\centering
\includegraphics[width=.98\textwidth]{nonlinear_marked_atlas.pdf}
\caption{The marked system, its complete packet jets, its operation-labeled tree grammar, and its exact packet equations are reversible charts.  The fixed point by itself is directional because it forgets the generating system.}
\label{fig:block11-marked-atlas}
\end{figure}

\subsection{Exact diagonal packet Fa\`a di Bruno compression}
\label{subsec:block11-faa-compression}

The packet Fa\`a di Bruno theorem from the nonlinear source expands an $n$th derivative over all set partitions and all intermediate packet labels \cite{P3,RT,NLO}.  When all directions are equal, the symmetric derivative tensors make many labeled terms identical.  They can be grouped exactly by block-size and packet-color multiplicities.

Let $F:Y\to Z$ and $G:X\to Y$ be $n$ times differentiable.  Fix an input packet $\beta$, an output packet $\delta$, a point $x$, and a direction $h\in P_\beta X$.  For every intermediate packet $\gamma$ and every $r\ge1$, put
\[
 g_{r,\gamma}
 =
 Q_\gamma D^rG(x)[h,\ldots,h].
\]
A colored partition profile is a finite family
\[
 \mu=(\mu_{r,\gamma})_{r\ge1,\,\gamma\in\Gamma}
\]
of nonnegative integers satisfying
\[
 \sum_{r,\gamma}r\mu_{r,\gamma}=n.
\]
Let
\[
 k(\mu)=\sum_{r,\gamma}\mu_{r,\gamma}.
\]

\begin{theorem}[Compressed diagonal packet Fa\`a di Bruno formula]
\label{thm:block11-compressed-faa}
With the notation above,
\[
\boxed{
\begin{aligned}
&P_\delta D^n(F\circ G)(x)[h,\ldots,h]\\
&=\sum_{\mu}
\frac{n!}
{\displaystyle\prod_{r,\gamma}
 (r!)^{\mu_{r,\gamma}}\mu_{r,\gamma}!}
\,P_\delta D^{k(\mu)}F(Gx)
\bigl[
 (g_{r,\gamma})^{\mu_{r,\gamma}}
\bigr]_{r,\gamma}.
\end{aligned}}
\tag{11.5}
\]
Here the last line means that $g_{r,\gamma}$ is inserted $\mu_{r,\gamma}$ times; symmetry makes the insertion order irrelevant.
\end{theorem}

\begin{proof}
Start with the packet Fa\`a di Bruno sum over a set partition $\pi$ of $\{1,\ldots,n\}$ and a packet color attached to every block.  Because every input direction is $h$, the value of one term depends only on the number $\mu_{r,\gamma}$ of blocks of size $r$ carrying color $\gamma$.  The symmetric group acts transitively on all colored set partitions with one fixed profile.  The stabilizer permutes the $r$ entries inside each block and permutes the indistinguishable blocks of a fixed type $(r,\gamma)$.  Its order is
\[
 \prod_{r,\gamma}(r!)^{\mu_{r,\gamma}}\mu_{r,\gamma}!,
\]
so the orbit has the coefficient printed in \textup{(11.5)}.  Summing one representative per orbit gives the result.
\end{proof}

\begin{corollary}[Colored integer-partition term count]
\label{cor:block11-compressed-count}
If $Y$ has $M$ intermediate packets, the number of compressed diagonal profiles at order $n$ is
\[
 \boxed{
 C_{n,M}
 =[x^n]\prod_{r\ge1}(1-x^r)^{-M}.}
 \tag{11.6}
\]
Before compression, the number of packet-labeled set-partition terms is
\[
 N_{n,M}=\sum_{k=1}^{n}S(n,k)M^k.
\]
\end{corollary}

\begin{proof}
For each block size $r$, the multiplicities among $M$ colors contribute
\[
 (1-x^r)^{-M}.
\]
Multiplying over all $r$ and taking the coefficient of $x^n$ counts exactly the profiles satisfying the weighted-size constraint.  The raw count is the existing packet term-count formula, grouped by the number of blocks.
\end{proof}

\begin{table}[htbp]
\centering
\caption{Raw and compressed diagonal Fa\`a di Bruno term counts for three intermediate packets.}
\label{tab:block11-faa-counts}
\small
\begin{tabular}{rrrr}
\toprule
Order $n$ & Raw packet terms & Colored profiles & Compression factor\\
\midrule
1 & 3 & 3 & 1.0\\
2 & 12 & 9 & 1.3\\
3 & 57 & 22 & 2.6\\
4 & 309 & 51 & 6.1\\
5 & 1,866 & 108 & 17.3\\
6 & 12,351 & 221 & 55.9\\
7 & 88,563 & 429 & 206.4\\
8 & 681,870 & 810 & 841.8\\
9 & 5,597,643 & 1,479 & 3784.7\\
10 & 48,718,569 & 2,640 & 18454.0\\
\bottomrule
\end{tabular}
\end{table}


\begin{figure}[htbp]
\centering
\includegraphics[width=.88\textwidth]{faa_di_bruno_compression.pdf}
\caption{For repeated directions, packet-labeled set partitions collapse to colored integer-partition profiles.  The compression is exact, not asymptotic.}
\label{fig:block11-faa-compression}
\end{figure}

\begin{boundary}[Scope of the compression theorem]
Equation \textup{(11.5)} closes the repeated-direction, single-input-packet case.  For fully mixed directions, blocks carry the types of their actual input subsets; the relevant orbit type is finer and can approach the full set-partition expansion.  No universal mixed-direction compression ratio is asserted.
\end{boundary}

\subsection{The cyclic nonlinear core of a packet system}
\label{subsec:block11-dependency-core}

For the packet system \textup{(11.4)}, define the directed dependency graph $G_\Phi$ on $I$ by placing an edge
\[
 j\longrightarrow i
\]
when some active block contributing to the equation for $u_i$ contains $u_j$ as an input.  Let
\[
 \operatorname{SCC}(G_\Phi)
\]
be its strongly connected components.  A component is called cyclic when it contains more than one address or is a singleton carrying a self-loop.

\begin{definition}[Address-preserving cyclic dimension]
The cyclic packet dimension of the marked system is
\[
 \boxed{
 \operatorname{cdim}_{\mathcal P}(\Phi)
 =
 \sum_{C\text{ cyclic SCC}}|C|.}
 \tag{11.7}
\]
It measures the number of packet variables that remain in coupled nonlinear components after exact elimination of the acyclic singleton components.
\end{definition}

\begin{theorem}[Exact SCC reduction]
\label{thm:block11-scc-reduction}
Order the SCC condensation graph topologically.  Every acyclic singleton component can be substituted exactly once its predecessors have been determined.  Repeating this elimination gives a solution-preserving bijection between:
\begin{enumerate}
\item solutions of the full packet system, and
\item solutions of the subsystem supported on the cyclic SCCs, followed by deterministic reconstruction along the condensation graph.
\end{enumerate}
Consequently, the coupled nonlinear core has exactly $\operatorname{cdim}_{\mathcal P}(\Phi)$ packet variables.
\end{theorem}

\begin{proof}
A singleton SCC without a self-loop has an equation in which its own variable does not occur.  Once all predecessor components are known, that equation computes the variable directly.  Substituting this expression into every successor equation is an algebraic identity and gives a bijection of solution sets.  The condensation graph is acyclic, so repeated topological elimination terminates.  Components containing a directed cycle cannot be removed by this one-pass address-preserving substitution because each contains a variable depending, directly or transitively, on itself.
\end{proof}

\begin{figure}[htbp]
\centering
\includegraphics[width=.94\textwidth]{nonlinear_dependency_core.pdf}
\caption{Condensation separates exact one-pass reconstruction from the cyclic nonlinear core.  The resulting dimension is relative to the declared packet coordinates.}
\label{fig:block11-dependency-core}
\end{figure}

\begin{boundary}[Not a universal nonlinear state minimum]
The quantity \(\operatorname{cdim}_{\mathcal P}\) is invariant under relabeling of the declared packet system and exact elimination within its dependency grammar.  A nonlinear change of variables, hidden-variable introduction, algebraic elimination, or another packet family can change it.  It is therefore a certified address-preserving system dimension, not a universal minimum over all nonlinear presentations.
\end{boundary}

\subsection{A finite residue signature for simply generated trees}
\label{subsec:block11-tree-signature}

Consider the scalar marked equation
\[
 \boxed{
 u(z)=z+\Phi(u(z)),
 \qquad
 \Phi(y)=\sum_{m\in\mathcal A}w_my^m,}
 \tag{11.8}
\]
where $\mathcal A\subseteq\{2,3,\ldots\}$ is finite and every $w_m$ is a positive integer.  Its unique formal solution belongs to $z\mathbb Z[[z]]$ and counts plane rooted trees whose internal arities lie in $\mathcal A$, weighted by the $w_m$.

Define the arity period
\[
 \boxed{
 d_{\mathcal A}
 =
 \gcd\{m-1:m\in\mathcal A\}.}
 \tag{11.9}
\]
Every tree with internal vertices $v$ has leaf count
\[
 \ell(T)=1+\sum_{v}(\deg(v)-1).
\]

\begin{theorem}[Finite tree-address signature]
\label{thm:block11-tree-address}
Fix a packet width $N$ and put
\[
 g_N=\gcd(N,d_{\mathcal A}).
\]
For the arithmetic packets
\[
 U_j(x)=\sum_{r\ge0}[z^{Nr+j}]u(z)\,x^r,
\]
the following are equivalent:
\begin{enumerate}
\item $U_j\neq0$;
\item $j-1$ belongs to the subgroup of $\mathbb Z/N\mathbb Z$ generated by $\{m-1:m\in\mathcal A\}$;
\item
\[
 \boxed{j\equiv1\pmod{g_N}.}
\]
\end{enumerate}
Hence the exact number of active packet addresses is
\[
 \boxed{
 \#\operatorname{Supp}_N(u)=\frac{N}{g_N}.}
 \tag{11.10}
\]
\end{theorem}

\begin{proof}
Replacing one leaf by an $m$-ary vertex increases the leaf count by $m-1$.  Starting from the one-leaf tree, every nonnegative sum of allowed increments is realized by a sequence of leaf replacements.  Positivity of the weights prevents cancellation, so the coefficient support is exactly
\[
 1+\left\langle m-1:m\in\mathcal A\right\rangle_{\mathbb N}.
\]
Modulo $N$, the monoid generated by the increments equals the subgroup they generate because every element of a finite cyclic subgroup has a nonnegative representative.  That subgroup is $g_N\mathbb Z/N\mathbb Z$, which proves the congruence and the count.
\end{proof}

\begin{corollary}[Finite bottom-up residue automaton]
Label every leaf by $1\in\mathbb Z/N\mathbb Z$ and every $m$-ary vertex by the sum of its child labels.  The reachable state set is exactly
\[
 \{j:j\equiv1\pmod{g_N}\}.
\]
Thus the finite pair
\[
 \boxed{
 \mathfrak T_N(\mathcal A)
 =
 (d_{\mathcal A},g_N,N/g_N)}
 \tag{11.11}
\]
classifies packet support for the marked simply generated family.
\end{corollary}

\begin{figure}[htbp]
\centering
\includegraphics[width=.72\textwidth]{tree_residue_signature.pdf}
\caption{For arities three and five, the increment gcd is two.  At width twelve, exactly the six addresses congruent to one modulo two are active.}
\label{fig:block11-tree-residue}
\end{figure}

\subsection{Dominant singularity orbits}
\label{subsec:block11-tree-singularity}

Assume now that the coefficients $w_m$ are positive real numbers.  The classical smooth implicit-function schema for simply generated trees applies when there is a positive solution $\tau$ of
\[
 \Phi'(\tau)=1
\]
and
\[
 \rho=\tau-\Phi(\tau)>0.
\]
Under its standard nondegeneracy hypotheses, the reduced aperiodic series has a square-root singularity at its positive radius \cite{FlajoletSedgewick2009}.  The arity period restores the complete source orbit.

\begin{theorem}[Complete dominant-orbit signature]
\label{thm:block11-dominant-orbit}
Under the smooth implicit schema, the dominant singularities of $u(z)$ form
\[
 \boxed{
 \rho\,\mu_{d_{\mathcal A}},}
 \tag{11.12}
\]
where $\mu_d$ is the group of $d$th roots of unity.  After packetization with $x=z^N$, the union of dominant singularities of the complete packet tuple is
\[
 \boxed{
 \rho^N\,\mu_{h_N},
 \qquad
 h_N=\frac{d_{\mathcal A}}{\gcd(N,d_{\mathcal A})}.}
 \tag{11.13}
\]
Therefore the same finite arity signature determines both the number of active packet addresses and the size of the complete dominant-singularity orbit.
\end{theorem}

\begin{proof}
By \cref{thm:block11-tree-address}, the coefficient support is contained in $1+d_{\mathcal A}\mathbb N$ and the reduced exponent set is aperiodic by maximality of the gcd.  Hence
\[
 u(z)=zV(z^{d_{\mathcal A}})
\]
with $V$ aperiodic.  The smooth implicit schema gives one positive dominant singularity for $V$, and pulling back through $z^{d_{\mathcal A}}$ gives \textup{(11.12)}.  The map $z\mapsto z^N$ sends $\mu_{d_{\mathcal A}}$ onto a root group of cardinality
\[
 d_{\mathcal A}/\gcd(N,d_{\mathcal A}).
\]
No packet can introduce a new dominant image outside this set.  Conversely, if every packet were analytic at one image, the exact synthesis
\[
 u(z)=\sum_{j=0}^{N-1}z^jU_j(z^N)
\]
would make $u$ analytic at every source point above that image, contradicting \textup{(11.12)}.  Thus the union of packet singularities is exactly \textup{(11.13)}.
\end{proof}

\begin{figure}[htbp]
\centering
\includegraphics[width=.94\textwidth]{tree_singularity_orbit.pdf}
\caption{Arithmetic packetization maps the source root-of-unity singularity orbit through $z\mapsto z^N$.  Distinct source singularities can alias to one packet-variable singularity.}
\label{fig:block11-tree-singularity}
\end{figure}

\begin{boundary}[Individual-packet cancellation]
The theorem identifies the union of dominant singularities of the complete packet tuple.  One individual packet can lose some orbit points through address cancellation.  The finite signature does not, by itself, determine every individual local amplitude or cancellation coefficient.
\end{boundary}

\subsection{\texorpdfstring{The sharp full \(q\)-ary family}{The sharp full q-ary family}}
\label{subsec:block11-qary}

Take
\[
 \boxed{u=z+u^q,\qquad q\ge2.}
 \tag{11.14}
\]
Lagrange inversion gives the exact coefficient formula
\[
 \boxed{
 [z^n]u
 =
 \begin{cases}
 \displaystyle
 \frac{1}{(q-1)m+1}\binom{qm}{m},
 &n=1+(q-1)m,\\[9pt]
 0,&n\not\equiv1\pmod{q-1}.
 \end{cases}}
 \tag{11.15}
\]
The nonzero coefficients are the Fuss--Catalan numbers of full plane $q$-ary trees.

Write
\[
 u(z)=\sum_{j=0}^{N-1}z^jU_j(z^N),
 \qquad x=z^N.
\]

\begin{theorem}[Exact $q$-ary packet phase theorem]
\label{thm:block11-qary-phase}
Put
\[
 g=\gcd(N,q-1),
 \qquad
 h=\frac{q-1}{g}.
\]
Then:
\begin{enumerate}
\item the active packet addresses are exactly
\[
 \boxed{j\equiv1\pmod g;}
\]
\item the complete packet system is
\[
\boxed{
 U_j(x)
 =
 \delta_{j,1}
 +
 \sum_{\substack{0\le r_1,\ldots,r_q<N\\
 r_1+\cdots+r_q\equiv j\ (\mathrm{mod}\ N)}}
 x^{(r_1+\cdots+r_q-j)/N}
 \prod_{a=1}^{q}U_{r_a}(x);}
 \tag{11.16}
\]
\item the smooth critical values are
\[
 \tau=q^{-1/(q-1)},
 \qquad
 \rho=(q-1)q^{-q/(q-1)};
\]
\item the source dominant orbit has $q-1$ points, while the complete packet orbit has exactly $h$ points.
\end{enumerate}
The phase is therefore determined by the finite tuple
\[
 \boxed{(q,N,g,N/g,h).}
 \tag{11.17}
\]
\end{theorem}

\begin{proof}
Formula \textup{(11.15)} follows from Lagrange inversion applied to
\[
 z=u(1-u^{q-1}).
\]
The support and active-address statements are \cref{thm:block11-tree-address} with $d_{\mathcal A}=q-1$.  Expanding $u^q$ in complete packets and collecting the exponent carry
\[
 r_1+\cdots+r_q=j+Nc
\]
gives \textup{(11.16)}.  The critical equations
\[
 q\tau^{q-1}=1,
 \qquad
 \rho=\tau-\tau^q
\]
give the printed values, and \cref{thm:block11-dominant-orbit} gives the orbit count.
\end{proof}

\begin{table}[htbp]
\centering
\caption{Exact packet signatures for the marked full $q$-ary tree family $u=z+u^q$.}
\label{tab:block11-qary-phases}
\small
\begin{tabular}{rrrrrr}
\toprule
$q$ & $N$ & $g=\gcd(N,q-1)$ & Active packets & Packet orbit & Active addresses\\
\midrule
2 & 5 & 1 & 5 & 1 & $\{0,1,2,3,4\}$\\
3 & 4 & 2 & 2 & 1 & $\{1,3\}$\\
3 & 5 & 1 & 5 & 2 & $\{0,1,2,3,4\}$\\
4 & 4 & 1 & 4 & 3 & $\{0,1,2,3\}$\\
4 & 6 & 3 & 2 & 1 & $\{1,4\}$\\
5 & 6 & 2 & 3 & 2 & $\{1,3,5\}$\\
5 & 7 & 1 & 7 & 4 & $\{0,1,2,3,4,5,6\}$\\
6 & 10 & 5 & 2 & 1 & $\{1,6\}$\\
\bottomrule
\end{tabular}
\end{table}


\subsection{Verification, status, and exact boundary}
\label{subsec:block11-validation}

The Nonlinear tree atlas verifier performs four independent exact audits.  It enumerates colored integer-partition profiles and compares their orbit-weighted Fa\`a di Bruno values with the full set-partition expansion.  It constructs exact fixed-point prefixes for marked polynomial tree families and checks the finite arity-gcd support theorem over many packet widths.  It verifies Fuss--Catalan formulas and the complete carry-form packet equations for the full $q$-ary family.  Finally, it checks dependency-graph reductions against exhaustive bounded graph instances.

\begin{table}[htbp]
\centering
\caption{Exact finite verification for Nonlinear tree atlas.}
\label{tab:block11-validation}
\small
\begin{tabularx}{\textwidth}{@{}Y r@{}}
\toprule
Validation class & Exact checks\\
\midrule
Colored profile counts & 100\\
Orbit-size identity & 48\\
Exact compressed Faà di Bruno values & 126\\
Polynomial fixed-point prefixes & 300\\
Finite tree-address signatures & 70,458\\
Fuss-Catalan formulas & 2,892\\
Exact q-ary packet systems & 37,567\\
Dominant-orbit combinatorics & 176\\
Dependency-core reductions & 4,054\\
\midrule
\textbf{Nonlinear tree atlas total} & \textbf{115,721}\\
Cumulative the preceding certified atlas program total & 2,027,665\\
Failures & 0\\
\bottomrule
\end{tabularx}
\end{table}


\begin{resultbox}
\textbf{Section result.}
Marked polynomial fixed-point systems admit exact packet, hypergraph, operadic-tree, and packet-equation charts.  The repeated-direction packet Fa\`a di Bruno expansion compresses exactly from colored set partitions to colored integer-partition profiles.  The packet dependency graph has an exact SCC reduction whose cyclic vertices form the address-preserving nonlinear core.  For positive simply generated tree systems, the arity gcd is a finite signature determining every arithmetic packet-support pattern and the complete dominant-singularity orbit.  The full $q$-ary family supplies a closed Fuss--Catalan phase theorem and exact packet equations.
\end{resultbox}

\begin{boundarybox}
\textbf{Exact boundary of Nonlinear tree atlas.}
The atlas classifies marked polynomial systems, not semantic fixed points independently of their generators.  The compression theorem treats repeated directions; arbitrary mixed-direction jets can require a finer orbit type.  The SCC dimension is relative to the declared packet coordinates.  The tree signature classifies support and the complete dominant-singularity orbit for the marked simply generated family.  It does not classify algebraic degree, universal minimal system dimension, every individual packet amplitude, arbitrary sign-canceling weights, noncontractive global dynamics, or all nonlinear operadic systems.
\end{boundarybox}
